\documentclass[10pt,a4paper]{amsart}
\usepackage[margin=19mm]{geometry}
\usepackage{amsmath,amssymb,mathtools}
\usepackage[scr=rsfs,cal=euler]{mathalfa}
\usepackage{xfrac}
\usepackage[unicode,hidelinks,bookmarks=false]{hyperref}
\newcommand{\ie}{i.e.\ }
\newcommand{\cf}{cf.\ }
\newcommand{\resp}{respectively\ }
\newcommand{\Subsection}{\subsection}
\providecommand{\nobreakdash}{\nobreak\hbox{-}}
\setlength{\emergencystretch}{2em}
\multlinegap0pt
\usepackage{mathtools}
\usepackage{amssymb}
\usepackage{xcolor}
\newtheorem{definition}{Definition}
\newtheorem{lemma}{Lemma}
\newtheorem{proposition}{Proposition}
\newtheorem{corollary}{Corollary}
\title{On intersection cohomology with torus \\ action of complexity one, II}
\author{Marta Agustin Vicente, Narasimha Chary Bonala, Kevin Langlois}
\date{Source: arXiv:2109.12641v2 (CC BY 4.0)}
\begin{document}
\maketitle

\noindent\emph{Source and licence.} On intersection cohomology with torus \\ action of complexity one, II. Version arXiv:2109.12641v2 (CC BY 4.0), distributed by arXiv under the Creative Commons Attribution 4.0 International licence, http://creativecommons.org/licenses/by/4.0/, as recorded in the arXiv licence metadata for this exact version and shown on the abstract page https://arxiv.org/abs/2109.12641v2. Full licence notice: \texttt{licenses/arxiv-2109.12641-CC-BY-4.0.txt}.\par\smallskip

\noindent\textit{Editorial scope.} Counted unit: the article's corollary xoro-expliteqvarT (source0.tex lines 1366--1369). The article's complete original proof of it is delivered with it (source0.tex lines 1370--1382, one \texttt{proof} environment of 13 lines). No mathematics is written, repaired or re-derived: the statement and the proof are the article's own lines, in the article's own notation, and no retained line is substituted or re-flowed.  Retained uncounted as the source-local context the proof needs: lines 800--816; lines 875--878; lines 893--896; lines 1306--1320; lines 1347--1355.  Nothing was omitted from any retained span, so the package carries no omission record.  No retained line contains a citation, so the wrapper carries no bibliography and no bibliography entry.  No retained line contains a figure environment, an image-inclusion command or drawing code, so no image is delivered and the package declares no asset dependency.  Editorial formatting only, in addition: an AMS \texttt{amsart} preamble that restates the article's own declarations and macros used by the retained lines, namely the environments \texttt{definition}, \texttt{lemma}, \texttt{proposition}, \texttt{corollary} on separate counters, together with the standard packages those lines need.  The retained environments print in this document as Definition 1, Lemma 1, Proposition 1, Corollary 1 in order of appearance, whereas the article prints the same items with the numbering of its own full document.  The complete native source of the article as distributed in the arXiv e-print for this exact version is bundled unchanged as \texttt{sources/116474/source0.tex}.
\begin{definition}\label{def-weightpack}
The \emph{weight matrix} of  $X$ is the matrix
${\rm Mat}_{E, B}(F)\in {\rm Mat}_{\ell\times n}(\mathbb{C})$
for the linear map $F: N\rightarrow \overline{N}$ corresponding to the inclusion $\mathbb{T}\hookrightarrow \mathbb{G}$. The \emph{matrix factorization} is the short exact sequence 
$$0\rightarrow N\xrightarrow{F} \overline{N}\xrightarrow{P} {\rm Coker}(F) =  \overline{N}/F(N)\rightarrow 0 $$
together with a splitting $S:\overline{N}\rightarrow N$, that is, $S$ is linear and $S\circ F = {\rm id}_{N}$ . Note that $F(N)$ must be a satured, i.e.
$$\overline{N}\cap F(N)_{\mathbb{Q}}\subset \overline{N}_{\mathbb{Q}}$$ 
is equal to $F(N)$.
Moreover, the map
\begin{equation}\label{Eq:downgraded}
    \phi:\overline{N}\rightarrow N\oplus {\rm Coker}(F),\,\, v\mapsto (S(v), P(v))
\end{equation}
is a $\mathbb{Z}$-module isomorphism. The \emph{quotient fan} $\Sigma$ of $X$ is the coarsest fan of ${\rm Coker}(F)_{\mathbb{Q}}$ containing all the strictly convex polyhedral cones $P(\delta_0)$, where $\delta_{0}$ runs over the face of the 
first quadrant 
$$ \delta:= \left\{ \sum_{i = 1}^{\ell}\alpha_{i} e_{i}\, \, |\, \alpha_{i}\in \mathbb{Q}_{\geq 0}, \,\, 1\leq i\leq \ell\right\}\subset \overline{N}_{\mathbb{Q}}.$$
The \emph{weight package} of $X$ is the data $\theta =  (\overline{N}, N, F, S, \Sigma)$.
\end{definition}

\begin{lemma}\label{lem-Ctheta-Ctheta}
With the notation of Definition \ref{def-weightpack}, the rational quotient $\psi: \mathbb{A}_{\mathbb{C}}^{\ell}\dashrightarrow X_{\Sigma}$
for the $\mathbb{T}$-action is regular on the open $\mathbb{G}$-orbit. Furthermore, $\widehat{C_{\theta}}$ is the normalization of the Zariski closure $C_{\theta}$ of $\psi(X\cap \mathbb{G})$ in $X_{\Sigma}$. 
\end{lemma}

\begin{definition}\label{def-pullbackpol}
The \emph{pullback polyhedral divisor of $(X, \theta)$} is 
$\bar{\mathfrak{D}}_{\theta}:= \kappa^{\star}(\mathfrak{D}_{\theta})$, where $\kappa: \widehat{C_{\theta}}\rightarrow X_{\Sigma}$ is the composition of the inclusion $C_{\theta}\subset X_{\Sigma}$ (see Lemma \ref{lem-Ctheta-Ctheta}) and the normalization. 
\end{definition}

\begin{lemma}\label{lem-equa-tvar}
Let $\theta =  (\overline{N}=  \mathbb{Z}^{\ell}, N  =  \mathbb{Z}^{n}, F, S,\Sigma)$ be a weight package.
Let $C\subset \mathbb{T}_{{\rm Coker}(F)}$ be an irreducible curve with defining ideal $I_{C} =  (f_{i}(x_{1}, \ldots, x_{s})\,|\, 1\leq i\leq a)$. Let $C_{\theta}$ be Zariski closure of $C$ in $X_{\Sigma}$. Consider the matrix 
$$P =   \begin{bmatrix} 
     b_{1,1} & \dots & b_{1, \ell}\\
    \vdots &     & \vdots  \\
     b_{s,1} & \dots & b_{s, \ell}
    \end{bmatrix}\in {\rm Mat}_{s\times \ell}(\mathbb{Z})$$
coming from the matrix factorization of $\theta$.
Set 
$$g_{i}(T_{1}, \ldots, T_{\ell}):= f_{i}\left(\prod_{j = 1}^{\ell}T_{j}^{b_{1, j}}, \ldots, \prod_{j = 1}^{\ell}T_{j}^{b_{s, j}}\right)\text{ for } 1\leq i\leq a,$$
and assume that there are Laurent monomials $u_{i}\in \mathbb{C}[T_{1}, T_{1}^{-1}, \ldots, T_{\ell}, T_{\ell}^{-1}]$ such
that $$X: = \mathbb{V}(u_{1}g_{1}, \ldots, u_{a}g_{a})\subset \mathbb{A}^{\ell}_{\mathbb{C}}$$ is irreducible. 
Then $X\subset \mathbb{A}^{\ell}_{\mathbb{C}}$ is a subvariety with linear torus action of complexity one and its normalization is equivariantly isomorphic to $X(\widehat{C_{\theta}}, \bar{\mathfrak{D}}_{\theta})$, where $\widehat{C_{\theta}}$ is the normalization of $C_{\theta}$ and $\bar{\mathfrak{D}}_{\theta}$ is pullback polyhedral divisor as in Definition \ref{def-pullbackpol}.
\end{lemma}

\begin{proposition}\label{prop-Construction2-irredhyper}
Let
$\theta =   (\overline{N}=  \mathbb{Z}^{\ell}, N  =  \mathbb{Z}^{n}, F, S,\Sigma)$ be a weight package. Consider
a Laurent polynomial of the form
$$g_{1}(T_{1}, \ldots, T_{\ell}):= f_{1}\left(\prod_{j = 1}^{\ell}T_{j}^{b_{1, j}}, \ldots, \prod_{j = 1}^{\ell}T_{j}^{b_{s, j}}\right),$$
where  $f_{1}\in\mathbb{C}[\mathbb{T}_{{\rm Coker}(F)}] =  \mathbb{C}[x_{1}, x_{1}^{-1}, \ldots, x_{s}, x_{s}^{-1}]$ is irreducible and the $b_{i, j}$ are the coefficients of the $P$-matrix of $\theta$. 
 Then there is a Laurent monomial $u_{1}\in \mathbb{C}[T_{1}, T_{1}^{-1}, \ldots, T_{\ell}, T_{\ell}^{-1}]$
such that $u_{1} g_{1}\in \mathbb{C}[T_{1}, \ldots, T_{\ell}]$ is irreducible. 
\end{proposition}

\begin{corollary}\label{xoro-expliteqvarT}
Any subvariety $X\subset \mathbb{A}^{\ell}_{\mathbb{C}}$ with linear torus action of complexity one has a
decomposition $X  =  \mathbb{V}(u_{1}g_{1}, \ldots, u_{a}g_{a})$ as in Lemma \ref{lem-equa-tvar}.
\end{corollary}

\begin{proof}
The $\mathbb{T}_{N}$-action on $\mathbb{A}^{\ell}_{\mathbb{C}}$ induces an $M$-grading on the ring $R:= \mathbb{C}[x_{1}, \ldots, x_{\ell}]$. Declare an element of $R$ homogeneous if it is homogeneous with respect to the $M$-grading.
The vanishing ideal of $X$ in $R$ is generated by non-constant homogeneous polynomials $h_{1}, \ldots, h_{a}\in R$. Set $X_{i}:= \mathbb{V}(h_{i})\subset \mathbb{A}^{\ell}_{\mathbb{C}}$ and consider the decomposition
$$ X_{i} =  \bigcup_{j =  1}^{s_{i}}X_{i, j}, \,\, 1\leq i\leq a$$
in irreducible components. Since $X_{i}$ is $\mathbb{T}_{N}$-stable, each component $X_{i, j}$ is of the form $\mathbb{V}(h_{i, j})\subset \mathbb{A}^{\ell}_{\mathbb{C}}$ for some homogeneous irreducible $h_{i, j}\in R$. Let $E$ be the set of pairs $(i, j)$ with $1\leq i\leq a$ and $1\leq j\leq s_{i}$
such that $h_{i, j}$ is not a scalar multiplication of a coordinate $x_{k}$. By 
Proposition \ref{prop-Construction2-irredhyper} we have  for $(i, j)\in E$ a decomposition $X_{i, j}  =  \mathbb{V}(u_{i, j} g_{i, j})$
obtained by pulling back $\psi(\mathbb{G}\cap X_{i, j}) = \mathbb{V}(f_{i, j})\subset \mathbb{T}_{{\rm Coker}(F)}$ by the quotient $\psi: \mathbb{G}\rightarrow \mathbb{T}_{{\rm Coker}(F)}$, where $u_{i, j}$ is a monomial. 
For $(i, j)\not\in E$ with $1\leq i\leq a$ and $1\leq j\leq s_{i}$ we set $u_{i, j} =  h_{i, j}$ and $g_{i, j} =  1$. Furthermore, for $1\leq i\leq a$
write 
$$ u_{i} =  \prod_{j = 1}^{s_{i}} u_{i, j},\,\, g_{i} =  \prod_{j =  1}^{s_{i}}g_{i, j} \text{ and }f_{i}  =  \prod_{j =  1}^{s_{i}}f_{i, j}.$$
Then $X  =  \mathbb{V}(u_{1}g_{1}, \ldots, u_{r}g_{a})$. Moreover, the zero locus $\mathbb{V}(f_{1}, \ldots, f_{r})\subset \mathbb{T}_{{\rm Coker}(F)}$ is the image $\psi(X\cap \mathbb{G})$, whence the result. 
\end{proof}

\end{document}
